\documentclass[11pt]{ctexart}
\usepackage[a4paper,top=2cm,bottom=2cm,left=2.2cm,right=2.2cm]{geometry}
\usepackage{fancyvrb}
\usepackage{xcolor}
\usepackage{colortbl}
\usepackage{fontspec}
\usepackage[most]{tcolorbox}
\usepackage{booktabs}
\usepackage{enumitem}
\usepackage{titlesec}
\usepackage{hyperref}

\setmonofont{Consolas}
\setCJKmonofont{SimHei}
\setlength{\parindent}{0pt}
\setlist{nosep,leftmargin=1.4em}

% ── Step2 主题色（与 Step1 深蓝同族、换为深紫罗兰） ──
\definecolor{deep}{RGB}{58,32,96}        % 主色：深紫罗兰
\definecolor{accent}{RGB}{126,60,168}    % 强调：紫罗兰
\definecolor{keyline}{RGB}{90,130,200}   % 知识点：蓝
\definecolor{keybg}{RGB}{238,244,255}
\definecolor{warnline}{RGB}{200,120,60}  % 警告：橙
\definecolor{warnbg}{RGB}{255,243,235}
\definecolor{okline}{RGB}{26,116,88}     % 成果：绿
\definecolor{okbg}{RGB}{240,252,244}
\definecolor{tipline}{RGB}{140,110,180}  % 提示：紫
\definecolor{tipbg}{RGB}{246,242,252}
\definecolor{codebg}{RGB}{249,249,252}   % 代码底
\definecolor{huered}{RGB}{220,60,60}
\definecolor{huegrn}{RGB}{60,180,90}
\definecolor{hueblu}{RGB}{70,110,220}
\definecolor{rowtint}{RGB}{240,236,248}  % 表头浅紫

\titleformat{\section}{\large\bfseries\color{deep}}{\thesection}{0.6em}{}
\titleformat{\subsection}{\normalsize\bfseries\color{accent}}{\thesubsection}{0.5em}{}
\titleformat{\part}[display]{\centering\Huge\bfseries\color{deep}}{}{0pt}{\Huge}
\titlespacing{\section}{0pt}{14pt}{6pt}
\titlespacing{\subsection}{0pt}{10pt}{4pt}

\newtcolorbox{keybox}[1]{enhanced, colback=keybg, colframe=keyline,
  boxrule=0.6pt, arc=1.5mm, left=3mm, right=3mm, top=2mm, bottom=2mm,
  title={\bfseries #1}, fonttitle=\small\bfseries, coltitle=white,
  colbacktitle=keyline}
\newtcolorbox{warnbox}[1]{enhanced, colback=warnbg, colframe=warnline,
  boxrule=0.6pt, arc=1.5mm, left=3mm, right=3mm, top=2mm, bottom=2mm,
  title={\bfseries #1}, fonttitle=\small\bfseries, coltitle=white,
  colbacktitle=warnline}
\newtcolorbox{okbox}[1]{enhanced, colback=okbg, colframe=okline,
  boxrule=0.6pt, arc=1.5mm, left=3mm, right=3mm, top=2mm, bottom=2mm,
  title={\bfseries #1}, fonttitle=\small\bfseries, coltitle=white,
  colbacktitle=okline}
\newtcolorbox{tipbox}[1]{enhanced, colback=tipbg, colframe=tipline,
  boxrule=0.5pt, arc=1.5mm, left=3mm, right=3mm, top=1.5mm, bottom=1.5mm,
  title={\bfseries #1}, fonttitle=\small\bfseries, coltitle=white,
  colbacktitle=tipline}

\newcommand{\code}[1]{\texttt{#1}}
\DefineVerbatimEnvironment{Verilog}{Verbatim}{frame=single,fontsize=\small}

\begin{document}

% ==================== 封面 ====================
\begin{titlepage}
\centering
\vspace*{2.2cm}
{\color{deep}\Huge\bfseries 颜色识别与质心}\\[6pt]
{\color{deep}\Large\bfseries 从零手搓 Step2 教程}\\[14pt]
{\large 2026 FPGA 创新设计赛道（安路）赛题二}\\[6pt]
{\large 视觉方向 · Step2\_color\_centroid}\\[20pt]

\begin{tcolorbox}[enhanced, colback=codebg, colframe=accent, boxrule=0.8pt,
  arc=2mm, width=12.5cm, left=4mm, right=4mm, top=3mm, bottom=3mm]
\centering
{\small 读者：零基础队员（建议先读完 Step1 教程）\\
目标：理解并手搓视觉检测链——从像素流里"认出"目标并算出坐标\\
前置要求：会 Verilog 基础语法、跑过 ModelSim 仿真\\
预计学习时间：12～16 天（每天 2～3 小时）}
\end{tcolorbox}

\vspace{8mm}
\begin{tcolorbox}[enhanced, colback=tipbg, colframe=tipline, boxrule=0.5pt,
  arc=2mm, width=12.5cm, left=4mm, right=4mm, top=2mm, bottom=2mm]
\centering
{\small\color{deep} \textbf{本教程与你手上的代码是什么关系}\\
Step1 学会了"把像素接住"；Step2 要做的，是从这些像素里\textbf{认出颜色目标}\\
并算出它的\textbf{画面坐标}——这是自动跟踪的"眼睛"。}
\end{tcolorbox}

\vfill
{\small\color{black!50} fpga-ptz-tracker 教程系列 · 2026-09-19}
\end{titlepage}

\tableofcontents
\newpage

% ============================================================
\part{从零学习教程}

\section{你在造什么}

\subsection{一句话}

\textbf{从每秒 5500 万个像素里，实时找出"红色目标"在哪里，并给出坐标。}

上一份教程（Step1）你学会了"把摄像头的数据接住"；这一份要在这条像素流上做三件事：

\begin{keybox}{Step2 要回答的三个问题}
\begin{enumerate}
  \item \textbf{这个像素是什么颜色？}\quad 把 RGB 换算成"人更容易判定的颜色描述"（HSV）
  \item \textbf{它是不是我们要的目标色？}\quad 用一个"阈值窗"把画面上属于目标的像素挑出来（掩码）
  \item \textbf{目标在画面哪里？}\quad 把所有目标像素的坐标平均，得到"重心"（质心）；顺带算包围框
\end{enumerate}
\end{keybox}

\subsection{整条检测链}

\begin{Verbatim}[frame=single,fontsize=\small]
 96bit 像素流 ──> ser96to24 ──> rgb2hsv ──> hsv_threshold
 （Step1 已做）    1 像素/拍     RGB→HSV      挑出目标"红色"像素
                                                       │
                                                       ▼
 目标坐标 <── centroid_accum <── bin_smooth3x3 <── 掩码（1 像素 1 bit）
（cx, cy）     帧内累加求质心      3×3 多数表决去噪
\end{Verbatim}

整条链每拍处理 1 个像素，74MHz 时钟下每秒能过 5500 万像素——
和摄像头 720p60 的像素率完全匹配，不多不少。

\subsection{做完是什么效果}

屏幕上：红色目标上出现绿色十字准星和红色包围框，左上角状态行显示坐标：
\begin{Verbatim}[frame=single,fontsize=\small]
TRK0 LOCK X0560Y0320     ← 认出目标（LOCK），坐标 (560, 320)
TRK0 LOST  X----Y----    ← 画面里没有目标（LOST）
\end{Verbatim}

\begin{okbox}{学完 Step2 你能做到}
\begin{itemize}
  \item 解释 RGB 和 HSV 的区别，说清为什么判颜色要用 HSV
  \item 手算一个像素的 HSV 值和阈值判决结果
  \item 理解"查表代替除法"“移位代替乘除"这些 FPGA 常用技巧
  \item 看懂行缓存、3×3 窗口、质心累加器的电路结构
  \item 跑通检测链的 5 个测试，并知道每个测试在验什么
\end{itemize}
\end{okbox}

% ============================================================
\section{阶段 0：图像与颜色基础（1～2 天）}

\subsection{像素：三个数描述一个颜色}

屏幕上每个小格子叫一个\textbf{像素}（pixel），由三个数字组成：

\begin{Verbatim}[frame=single,fontsize=\small]
像素 = (R, G, B)    每项 0～255

(255,   0,   0) = 纯红        (  0, 255,   0) = 纯绿
(  0,   0, 255) = 纯蓝        (255, 255, 255) = 白
(128, 128, 128) = 中灰        (  0,   0,   0) = 黑
(224,  32,  32) = 暗一点的红  ← 本项目要跟踪的目标色
\end{Verbatim}

一个 1280×720 的画面有 92 万个这样的三元组，每秒刷新 60 次。
FPGA 要在 16.7 毫秒内把每一帧全部处理完——这就是"实时"的含义。

\subsection{RGB 判颜色的麻烦}

想找出"红色目标"，最直接的想法：

\begin{Verbatim}[frame=single,fontsize=\small]
if (R > 150) && (G < 80) && (B < 80) then 是红色
\end{Verbatim}

这个写法在实验室灯光下能工作，但一换环境就崩：

\begin{center}
\renewcommand{\arraystretch}{1.25}
\begin{tabular}{llll}
\toprule
\rowcolor{rowtint}\textbf{场景} & \textbf{同一块红布的颜色} & \textbf{R>150？} & \textbf{判定结果} \\
\midrule
室内正常光 & (224, 32, 32) & 是 & 正确识别 \\
室外强光 & (255, 120, 120) & 是 & 勉强正确（G/B 已超标） \\
阴影下 & (90, 12, 12) & 否 & \textcolor{huered}{漏检！} \\
白炽灯偏黄 & (240, 110, 60) & 是 & 勉强正确 \\
\bottomrule
\end{tabular}
\end{center}

\textbf{问题根源：RGB 三个数字里，"有多亮"和"是什么颜色"是混在一起的。}
光照变强时三个数一起变大；光照变弱时一起变小。
你想判断"是不是红色"，却不得不同时应付亮度的干扰。

\subsection{HSV：换个说法描述颜色}

把"三个混在一起的数"换成"三个各管一事的数"：

\begin{keybox}{HSV 三个分量}
\begin{itemize}
  \item \textbf{H（Hue，色相）}：这是\textbf{什么颜色}。0\textasciitilde 360 度一圈
  \item \textbf{S（Saturation，饱和度）}：颜色\textbf{有多纯}。0 = 灰，越大越鲜艳
  \item \textbf{V（Value，明度）}：颜色\textbf{有多亮}。0 = 黑，越大越亮
\end{itemize}
\end{keybox}

色相环（H 的取值）：

\begin{Verbatim}[frame=single,fontsize=\small]
                        0° 红
                 330° ──┐  ┌── 30°
                       │  │
            300° 品红 ──┤  ├── 60° 黄
                       │  │
            270°     ──┤  ├── 90°
                       │  │
            240° 蓝  ──┤  ├── 120° 绿
                       │  │
                 210° ──┘  └── 150°
                       180° 青

 判"红"只需：H 在 335°～360° 或 0°～25°（红色在环上是跨 0 度的一段）
 光照变化主要影响 S 和 V，对 H 的影响小得多——这就是 HSV 的好处
\end{Verbatim}

\begin{tipbox}{一句话记住}
RGB 是"三盏灯的亮度"；HSV 是"什么颜色 + 多纯 + 多亮"。
要和光照变化对抗，就用后者。
\end{tipbox}

\subsection{练习}

\begin{enumerate}
  \item 打开 Python，装好 \code{opencv-python}，运行：
\begin{Verbatim}[fontsize=\small]
import cv2, numpy as np
rgb = np.uint8([[[224, 32, 32]]])        # 一个红色像素（注意 cv2 要 BGR！）
hsv = cv2.cvtColor(np.uint8([[[32, 32, 224]]]), cv2.COLOR_BGR2HSV)
print(hsv)   # 看看输出的 H/S/V 是多少
\end{Verbatim}
  \item 找一张照片，看"红色物体"和"白色墙"的 HSV 分别是什么样
  \item 思考题：白色的 H 和 S 应该是多少？（提示：白色的"颜色"是什么？）
\end{enumerate}

% ============================================================
\section{阶段 1：定点算术——FPGA 的数学（2～3 天）}

\subsection{FPGA 不会做小数}

FPGA 里没有"浮点数"概念（至少本项目不用）。所有数都是整数。
那"除以 delta"这种运算怎么办？两个办法：\textbf{移位}和\textbf{查表}。

\begin{keybox}{技巧 1：移位代替乘除 2 的幂}
\begin{Verbatim}[fontsize=\small]
x << 1  等于  x × 2
x << 2  等于  x × 4
x >> 1  等于  x ÷ 2（向下取整）
x >> 6  等于  x ÷ 64
\end{Verbatim}
移位在电路里几乎不花资源（就是接线方式变了），硬件里是"免费"的。
\end{keybox}

\begin{keybox}{技巧 2：查表代替除法}
色相计算需要"除以 delta"（delta 是 1\textasciitilde 255 之间某个数）。
除法的电路很大很慢——但我们可以在\textbf{设计阶段}把 256 个可能的商全算出来，
预先放进一张小表（ROM），运行时"查一下"就行：

\begin{Verbatim}[fontsize=\small]
inv[1] = 3840, inv[2] = 1920, inv[3] = 1280, ...
inv[d] = floor(3840 / d)      d = 1～255

表大小：256 项 × 12 bit ≈ 3.8 Kbit（一块小 ROM 就装下）
运行时：查表只要 1 拍，比除法器又快又省
\end{Verbatim}
\end{keybox}

\subsection{截断陷阱：>> 是"向下取整"}

\begin{warnbox}{负数右移会变小}
\begin{Verilog}
-7 >> 1  =  -4     （不是 -3！向下取整：-3.5 → -4）
\end{Verilog}
本项目的色相计算里有负数（G−B 可能为负），右移的取整方向必须和
Python 金标准模型完全一致，否则 RTL 和模型会差 1，验证不通过。\\[4pt]
\textbf{规则}：Python 的 \code{>>} 和 Verilog 的带符号 \code{>>>} 都是向下取整——
两边写法对齐，就能逐位一致。
\end{warnbox}

\subsection{定点验证的重要性}

\begin{okbox}{本项目的工作方式：金标准模型}
\begin{enumerate}
  \item 先用 Python 写一个"标准答案"版本（\code{tools/golden\_model.py}）
  \item 再让 RTL 实现同样的算法，要求\textbf{逐位一致}（不是"差不多"）
  \item 用 1 万条随机向量测试，任何一位不同都算失败
\end{enumerate}
这叫"位精确（bit-exact）验证"。因为 FPGA 里没有"约等于"，
差 1 就可能让阈值判决翻转，目标时有时无。
\end{okbox}

\subsection{练习}

\begin{enumerate}
  \item 手算：\code{5 >> 1}、\code{-5 >> 1}、\code{7 >> 2}、\code{-7 >> 2}
  \item 用 Python 验证你的答案（注意 Python 的负数是无限精度，\code{-5 >> 1 = -3}——想想为什么和 Verilog 不同）
  \item 打开 \code{tools/golden\_model.py}，找到 \code{build\_inv\_lut()} 函数，看看倒数表怎么建的
\end{enumerate}

% ============================================================
\section{阶段 2：RGB→HSV 的电路原理（3～4 天）}

这一节是整个 Step2 的核心。跟着五步走，你就能手推任意像素的 HSV。

\subsection{第 1 步：找最大、最小、差值}

\begin{Verbatim}[frame=single,fontsize=\small]
输入：R, G, B（各 8 bit）

max   = max(R, G, B)      三个数里最大的
min   = min(R, G, B)      三个数里最小的
delta = max - min         最大最小之差（"颜色的纯不纯"）

例子：(224, 32, 32)  →  max=224, min=32,  delta=192
\end{Verbatim}

电路实现：三个数两两比较，找最大最小各需 2 次比较（比较器很便宜）。

\subsection{第 2 步：明度 V 和饱和度 S}

\begin{Verbatim}[frame=single,fontsize=\small]
V = max          有多亮（纯红的 V=224，说明它挺亮）
S = delta        有多纯（纯红 delta=192；灰色三个数相等 → delta=0）

注意：标准 HSV 里 S 要除以 max 归一化（S = delta/max × 255），
本项目为了省掉一次除法，直接用 delta 代替——效果够用（详见 Part II）。
\end{Verbatim}

\subsection{第 3 步：色相 H——判断在哪个扇区}

看哪个分量最大：

\begin{Verbatim}[frame=single,fontsize=\small]
R 最大 → 颜色在 0°～60° 附近（红～黄）
G 最大 → 颜色在 120°～180° 附近（绿～青）
B 最大 → 颜色在 240°～300° 附近（蓝～品红）

本项目简化版：三个扇区，基值 0° / 120° / 240°，
    优先级 R > G > B（两个相等时算前者）

例子：(224, 32, 32)：R 最大 → 基值 0°
\end{Verbatim}

\subsection{第 4 步：扇区内插值}

在扇区里还要细分"偏移多少度"，公式：

\begin{Verbatim}[frame=single,fontsize=\small]
R 最大时：  h = (G - B) × 3840 / delta / 64     ← 用到查表！
（其他扇区换分子：G 最大用 B-R，B 最大用 R-G）

为什么是 3840？  因为 3840/64 = 60，正好是一个扇区的宽度（60°）
除以 64 用 >>6 移位实现

例子：(224, 32, 32)：G-B = 0  →  h = 0  →  H = 0° + 0 = 0°（就是纯红，对上了）
\end{Verbatim}

\begin{keybox}{查表怎么用}
\code{inv[delta]} 查到 \code{3840/delta} 的预计算值（12 bit），然后：\\[3pt]
\code{h = (G-B) × inv[delta] >> 6}\\[3pt]
一次乘法 + 一次移位，没有任何除法——这就是定点技巧的威力。
\end{keybox}

\subsection{第 5 步：拼装完整公式}

\begin{Verbatim}[frame=single,fontsize=\small]
H = 扇区基值 + clamp(h, 0, 60)

clamp 的原因：h 由差值算出，理论上在 0～60 之间，
    但数值误差可能溢出边界，钳一下保险

最终输出：
    H：9 bit（0～359）    S：8 bit（=delta）    V：8 bit（=max）
\end{Verbatim}

\begin{warnbox}{已知简化的代价（必须知道）}
本项目用"3 扇区"简化，H 只能落在：
\textbf{0～60°（红黄区）、120～180°（绿青区）、240～300°（蓝紫区）}。
中间色（如品红 300°、黄绿 90°）会被并入相邻扇区。\\[4pt]
\textbf{实际影响}：品红色（R=B>G）会被判成 H=0——和红色混在一起。
对策：靠 S/V 门限过滤 + 演示时避免品红干扰物（详见 Part II 的实测场景）。\\[4pt]
完整 6 扇区版本要多一次比较和更多逻辑，本项目为省资源选了 3 扇区。
\end{warnbox}

\subsection{2 级流水线}

\begin{Verbatim}[frame=single,fontsize=\small]
第 1 拍：算 max/min/delta，选好扇区差值，连同 RGB 一起存下来
第 2 拍：查表（拿 delta 查 inv）→ 乘法 → 移位 → clamp → 出 H/S/V

为什么分 2 拍？ 组合逻辑太长（比较+查表+乘法叠在一起），
74MHz 时钟下时序过不去；拆成 2 拍每拍都轻松
为什么不拆更多？ 延迟只为了让电路跑得动，多了没意义
\end{Verbatim}

\subsection{练习}

\begin{enumerate}
  \item 手算这 4 个像素的 HSV（用 \code{inv} 表近似）：
        \code{(255,0,0)}、\code{(0,255,0)}、\code{(0,0,255)}、\code{(128,128,128)}
  \item 和 Python 的答案对一下（\code{tools/golden\_model.py} 里有 \code{rgb\_to\_hsv}）
  \item 想想：\code{(255, 255, 0)}（纯黄）的 H 是多少？为什么这个简化版本给不出 60°？
\end{enumerate}

% ============================================================
\section{阶段 3：阈值分割（2 天）}

\subsection{二值化：每个像素过一遍筛子}

有了 HSV，判断"是不是目标色"就是三个条件同时成立：

\begin{Verilog}
mask = (H 在窗口内) && (S >= s_lo) && (V 在 [v_lo, v_hi] 内)
\end{Verilog}

输出 \code{mask} 是 1 个 bit：1 = 目标像素，0 = 背景。
整帧过完得到一张"黑白图"（掩码）。

\subsection{H 窗口环绕——红色跨 0 度的坑}

\begin{warnbox}{红色的窗口是"两段"}
红色在色相环上跨 0°：335°～360° \textbf{和} 0°～25° 都算红色。\\[4pt]
不能用简单的 \code{(H >= 335) \&\& (H <= 25)}——这个条件永远不成立（335 > 25）！\\[4pt]
\textbf{正确写法}：判断 \code{h\_lo > h\_hi} 时改用"或"：
\begin{Verilog}
in_h = (h_lo > h_hi) ? ((H >= h_lo) || (H <= h_hi))
                     : ((H >= h_lo) && (H <= h_hi));
\end{Verilog}
\end{warnbox}

\subsection{S/V 门限各管什么}

\begin{center}
\renewcommand{\arraystretch}{1.25}
\begin{tabular}{lll}
\toprule
\rowcolor{rowtint}\textbf{条件} & \textbf{挡住什么} & \textbf{例子} \\
\midrule
S $\geq$ s\_lo & 灰、白、黑（不纯的颜色） & 白墙 S=0，被挡掉 \\
V $\geq$ v\_lo & 太暗的（阴影里的噪声） & 暗红阴影 V=40，被挡掉 \\
V $\leq$ v\_hi & （可选）过曝的 & 反光高亮区 \\
\bottomrule
\end{tabular}
\end{center}

\subsection{8 组预置：现场按键切换}

\begin{Verbatim}[frame=single,fontsize=\small]
组 0: 红   H在[335,25]  S>=80   V在[60,255]     ← 默认，跟踪红色目标
组 1: 绿   H在[110,145] S>=60   V在[40,255]
组 2: 蓝   H在[200,250] S>=60   V在[40,255]
组 3: 橙黄 / 组 4: 青 / 组 5: 品红 / 组 6: 白 / 组 7: 黑
\end{Verbatim}

现场换目标颜色（比如评委拿个绿色球），按一下按键切到组 1 就行——
不用改代码、不用重新烧板子。

\subsection{练习}

\begin{enumerate}
  \item 手算：\code{(224,32,32)} 在组 0 → mask = ?（H=0, S=192, V=224）
  \item 手算：\code{(128,128,128)} 在组 0 → mask = ?（想想 S）
  \item 打开 \code{hsv\_threshold.v}，找到 case 语句里的 8 组参数
\end{enumerate}

% ============================================================
\section{阶段 4：去噪——3×3 多数表决（3～4 天）}

\subsection{噪声从哪来}

掩码里会混进"假目标"：传感器噪点、边缘反光、压缩伪影……
它们通常是零散的孤立点。看这张示意（1 = 目标像素）：

\begin{Verbatim}[frame=single,fontsize=\small]
    有噪声的掩码                去噪后
    0 0 0 0 0 0 0               0 0 0 0 0 0 0
    0 1 1 1 0 0 0               0 1 1 1 0 0 0
    0 1 1 1 0 1 0   ← 噪点      0 1 1 1 0 0 0
    0 1 1 1 0 0 0               0 1 1 1 0 0 0
    0 0 0 0 0 0 0               0 0 0 0 0 0 0
\end{Verbatim}

\subsection{3×3 窗口 + 多数表决}

对每个像素，看它和周围 8 个邻居（共 9 个），数一数有几个是 1：

\begin{Verbatim}[frame=single,fontsize=\small]
    邻居窗口（以中心像素 c 为核心）
    ┌───┬───┬───┐
    │ a │ b │ d │
    ├───┼───┼───┤
    │ e │ c │ f │     9 个格子里 1 的个数 >= 5  →  输出 1
    ├───┼───┼───┤     否则                        →  输出 0
    │ g │ h │ i │
    └───┴───┴───┘
\end{Verbatim}

\textbf{为什么是"多数表决"而不是"平均数"？}\quad
9 个里有 5 个以上是 1 才保留——孤立噪点（周围都是 0）被干掉，
真实目标（成片都是 1）被保留。这就像投票：群众的眼睛是雪亮的。

\subsection{边界怎么办：复制}

画面最外圈像素没有完整邻居。本项目的策略：\textbf{复制边界}
（画面外的像素按"最近的画面内像素"算）。

\begin{Verbatim}[frame=single,fontsize=\small]
左上角 (0,0) 的实际窗口（复制边界后）：    画面外的位置一律复制最近像素
    ┌───┬───┬───┐
    │ c │ c │ d │
    ├───┼───┼───┤
    │ c │ c │ f │      这样 (0,0) 也能得到一个完整的 9 格窗口
    ├───┼───┼───┤
    │ g │ h │ i │
    └───┴───┴───┘
\end{Verbatim}

\subsection{行缓存：怎么"看到上一行"}

像素是一行一行流过来的。要组成 3×3 窗口，需要同时看到
"上一行"和"上上一行"——但那些像素已经流走了。怎么办？

\begin{keybox}{行缓存（line buffer）}
用一个能存\textbf{一整行}的存储器，把流过去的像素留下来：\\
\begin{itemize}
  \item 当前行：正在流的像素
  \item 行缓存 1：存上一行（1280 个 bit）
  \item 行缓存 2：存上上一行
\end{itemize}
三个来源凑齐，3×3 窗口就齐了。
\end{keybox}

\begin{warnbox}{！本项目最难的坑（Part II 有完整复盘）}
"怎么对齐三行"看起来简单，实际连踩 5 个坑（FIFO 收支差、移位方向、
窗口取列差一拍、延迟线比设计短一格……）才做对。\\[4pt]
最终方案：\textbf{索引行缓存}——读写用同一个列指针，
"读到上一行的同一列"从结构上就成立，不需要手工推相位。\\[4pt]
初学建议：先接受"最终版是对的"去看 Part II 的复盘，
比试图自己从零推快得多。
\end{warnbox}

\subsection{练习}

\begin{enumerate}
  \item 给这张小掩码（5×5）手工做一遍 3×3 多数表决（边界复制）：
\begin{Verbatim}[fontsize=\small]
0 0 0 0 0
0 1 1 1 0
0 1 0 1 0     ← 中心有个洞，表决后会被填上吗？
0 1 1 1 0
0 0 0 0 0
\end{Verbatim}
  \item 用 Python 对同一张图跑 \code{golden\_model.smooth3x3\_majority}，对比你的答案
\end{enumerate}

% ============================================================
\section{阶段 5：质心与包围框（3 天）}

\subsection{质心：所有目标像素的"重心"}

目标是一块由很多像素组成的区域。用哪一点代表它的位置？——\textbf{质心}
（所有目标像素坐标的平均值）：

\begin{Verbatim}[frame=single,fontsize=\small]
假如目标是这 3 个像素：(x, y) = (10, 4), (11, 4), (12, 4)

cx = (10 + 11 + 12) / 3 = 11        ← x 的平均
cy = ( 4 +  4 +  4) / 3 = 4         ← y 的平均
目标位置 = (11, 4)
\end{Verbatim}

一般公式（m 为掩码：1 = 目标像素）：

\begin{Verbatim}[frame=single,fontsize=\small]
对全画面的每个像素：
    count = Σ m                目标像素总数
    sum_x = Σ (x × m)          "x 的和"，m=1 时才累加
    sum_y = Σ (y × m)

帧末（一帧过完）算一次：
    cx = sum_x / count
    cy = sum_y / count
\end{Verbatim}

\subsection{为什么用质心，不用包围框中心}

\begin{center}
\renewcommand{\arraystretch}{1.25}
\begin{tabular}{lll}
\toprule
\rowcolor{rowtint}\textbf{方案} & \textbf{优点} & \textbf{缺点} \\
\midrule
包围框中心 & 简单（4 个寄存器） & 目标形状不规则时偏离视觉中心 \\
\textbf{质心（选定）} & 反映"质量分布"，更符合直觉 & 需要一次除法 \\
\bottomrule
\end{tabular}
\end{center}

本项目两个都算：质心用来画准星和控云台，包围框用来画红框。

\subsection{帧累加器的电路结构}

\begin{Verbatim}[frame=single,fontsize=\small]
每个像素到来时（1 拍内完成）：
    if (mask) begin
        count <= count + 1;          // 21 bit 计数器
        sum_x <= sum_x + x;          // 32 bit 累加器
        sum_y <= sum_y + y;
    end

  x 计数器：每拍 +1，行尾清零（0～1279）
  y 计数器：每行 +1，帧首清零（0～719）

帧末（第 719 行的行尾）：
    启动除法状态机：sum_x / count → cx
                    sum_y / count → cy
    除法大约 70 拍完成——一帧的消隐期有 7 万多拍，绰绰有余
\end{Verbatim}

\begin{keybox}{为什么除法要放到帧末？}
除法器是"大电路"，如果每个像素都做除法（每拍一次），
要么资源爆炸要么降频。\\[4pt]
而质心只需要\textbf{每帧一个结果}——把除法安排在帧末的消隐期
（画面不出像素的那段空档），慢慢地算，完全不影响像素流。
\end{keybox}

\subsection{除零保护与最小像素门限}

\begin{warnbox}{两个必须处理的边界情况}
\begin{enumerate}
  \item \textbf{count = 0}（画面里没有目标）：除法除零会得到垃圾值——
        必须判断"count 为 0 就直接输出无效"，不做除法
  \item \textbf{count 很小}（比如 5 个杂点）：质心会被噪点带偏——
        设一个最小门限（本项目 30 个像素），不够就当没目标
\end{enumerate}
\end{warnbox}

\subsection{包围框：4 个极值寄存器}

\begin{Verbatim}[frame=single,fontsize=\small]
if (mask) begin
    if (x < min_x) min_x <= x;      // 最左
    if (x > max_x) max_x <= x;      // 最右
    if (y < min_y) min_y <= y;      // 最上
    if (y > max_y) max_y <= y;      // 最下
end
  （另外需要一个 seen 标志：帧内第一次见到目标时把四个值都设成当前坐标）
\end{Verbatim}

\subsection{练习}

\begin{enumerate}
  \item 手工算：三个目标像素 (10,4), (11,4), (12,4) 的质心和包围框
  \item 思考：如果目标是 L 形（缺一角），质心会在图形的内部还是外部？
  \item 读 \code{centroid\_accum.v} 的除法状态机部分，理解"恢复法除法"
\end{enumerate}

% ============================================================
\section{阶段 6：跑通检测链测试（1～2 天）}

\subsection{一键回归}

\begin{Verbatim}[frame=single,fontsize=\small]
cd D:\fpga-ptz-tracker\tools
python run_all.py
\end{Verbatim}

检测链相关的 5 个测试：

\begin{center}
\renewcommand{\arraystretch}{1.25}
\begin{tabular}{lll}
\toprule
\rowcolor{rowtint}\textbf{测试} & \textbf{验证什么} & \textbf{规模} \\
\midrule
tb\_rgb2hsv & RGB→HSV 位精确 & 10010 条随机+角点向量 \\
tb\_hsv\_threshold & 阈值判决位精确 & 4000 条 × 8 组预置 \\
tb\_bin\_smooth & 3×3 平滑位精确 & 12288 像素 \\
tb\_centroid & 质心/包围框/除零 & 3 帧迷你几何 \\
tb\_detect\_chain & 全链端到端 & 3 帧真实场景（含干扰物） \\
\bottomrule
\end{tabular}
\end{center}

\subsection{怎么读测试结果}

\begin{Verbatim}[frame=single,fontsize=\small]
--- tb_rgb2hsv ...
    ==> PASS                          ← 10010 条向量全部逐位一致

--- tb_detect_chain ...
    # FRAME0 OK v=1 cnt=8096 cx=604 cy=345
    # FRAME1 OK v=1 cnt=32592 cx=644 cy=391
    # FRAME2 OK v=1 cnt=221 cx=107 cy=108
    ==> PASS                          ← 3 帧场景全部与金标准一致
\end{Verbatim}

每个 \code{OK} 行显示的是检测结果：有效性、像素数、质心坐标。
和 \code{chain\_exp.txt}（Python 金标准算的期望值）逐字段对比。

\begin{okbox}{Step2 学习清单（全部打勾就过关）}
\begin{itemize}
  \item 能解释 RGB 和 HSV 的区别，说清 S/V 门限各挡什么
  \item 能手算一个像素的 HSV 和阈值判决
  \item 能解释"查表代替除法"和"移位代替乘除"
  \item 能画出行缓存 + 3×3 窗口的框图
  \item 能解释质心公式和除零保护
  \item 跑通 5 个测试，能说清每个测试在验什么
\end{itemize}
\end{okbox}

% ============================================================
% Step2 教程 · Part II：工程实现详解
% 由 step2_tutorial.tex \input 引入（共用同一份导言区与主题色）
% ============================================================

\part{工程实现详解}
\label{part:impl}

\begin{tipbox}{这一部分写什么}
上面的 Part I 是"从零学知识"；这一部分换一个角度：把仓库里\textbf{真实做出来的
Step2 工程}逐模块讲透——每个文件为什么这么写、每一步踩过什么坑、
怎么用金标准证明它是对的。学完 Part I 再读这里，你会看到"教科书知识"
是怎样变成仓库里 916 行可综合 RTL 的。

对应目录：\code{03\_vision/step2\_color\_centroid/}\\
对应代码：\code{03\_project/rtl/} 下 6 个 RTL 文件（916 行）\\
对应证据：\code{03\_project/sim/} 下的向量与期望文件、\code{04\_迭代文件/迭代记录\_v1到v7.md}
\end{tipbox}

% ============================================================
\section{工程总览}

\subsection{文件清单}

\begin{center}
\begin{tabular}{llr l}
\toprule
\rowcolor{rowtint}\textbf{文件} & \textbf{类型} & \textbf{行数} & \textbf{职责} \\
\midrule
\code{rgb2hsv.v} & RTL & 100 & RGB888 到 HSV 转换，2 级流水 \\
\code{hsv\_inv\_lut.v} & 生成 & 267 & 倒数表 ROM，脚本生成禁止手改 \\
\code{hsv\_threshold.v} & RTL & 58 & 8 组预置阈值二值化 \\
\code{bin\_smooth3x3.v} & RTL & 120 & 3x3 多数表决去噪（v7 索引行缓存） \\
\code{centroid\_accum.v} & RTL & 208 & 帧累加求质心/包围框 + 帧末除法 \\
\code{detect\_top.v} & RTL & 163 & 检测链顶层（96bit 窃流旁路） \\
\midrule
\code{tb/tb\_*.v} & 测试 & 5 个 & 单模块 4 个 + 整链 1 个 \\
\code{sim/*.hex/.txt} & 向量 & --- & 金标准生成的激励与期望 \\
\bottomrule
\end{tabular}
\end{center}

\subsection{数据流与延迟账}

\begin{Verilog}
96bit AXI-S (4px/beat, AWB 后)
   |  ser96to24 (Step1 模块, FIFO 降带宽)
   v
1 像素/拍 (24bit)
   |  rgb2hsv          延迟 2 拍
   v
H/S/V (9+8+8 bit)
   |  hsv_threshold    延迟 1 拍
   v
mask (1bit)
   |  bin_smooth3x3    延迟 ~3 拍 + 2 行行缓存
   v
mask_smooth
   |  centroid_accum   帧内累加, 帧末 ~66 拍除法(消隐期)
   v
目标包 (坐标/包围框/计数) -> 云台 / OSD
\end{Verilog}

整链延迟约 9 拍（不含帧累加），全部工作在 1 像素/拍 @ 74.25 MHz。
\textbf{关键约束}：除法必须在帧消隐期（约 1440 拍）内完成，66 拍绰绰有余，
所以选了最省资源的串行恢复法除法，而不是并行除法器。

\subsection{三条设计原则}

\begin{keybox}{本工程反复验证过的三条原则}
\textbf{1. 金标准先行}：所有算术先在 \code{tools/golden\_model.py} 用 Python
定点实现，RTL 逐位复刻。失配时先对照 Python 找差异，再改 RTL。\\
\textbf{2. 自对齐结构优先}：宁可多用一点逻辑资源，也要选"结构上不可能错位"的
方案（详见 bin\_smooth3x3 的 v3 到 v7 重构）。\\
\textbf{3. 流水线标志全程随行}：valid/sof/eol 三级标志跟着数据一起打拍，
任何模块都不许"自己数拍子猜边界"。
\end{keybox}

% ============================================================
\section{模块详解：rgb2hsv}

\subsection{为什么用 3840 这个数}

HSV 标准公式里有个除以 delta 的运算，FPGA 不做除法。做法是把
1/delta 提前算成表：\code{inv[d] = floor(3840/d)}。3840 的来历：

$$3840 = 60 \times 64$$

- 乘 60：让扇区内插值直接出"度"（每扇区 60 度）；
- 乘 64：最后用 \code{>>6} 除掉它，移位代替除法。

于是 \code{h = (diff × inv[delta]) >> 6} 的值域恰好是 0..60，
加扇区基值（0/120/240）就是完整色相。

\subsection{两级流水的分工}

\begin{Verilog}
第 0 拍（组合 + 打拍）:
  max/min 比较树, delta = max - min
  扇区判定: is_r = (R>=G && R>=B);  is_g = !is_r && (G>=R && G>=B)
  扇区差值: R 扇区取 (G-B), G 扇区取 (B-R), B 扇区取 (R-G)   [9bit 有符号]

第 1 拍（查表 + 乘 + 移位）:
  inv = inv_lut[delta]                     // 256x12 ROM
  h   = clamp((diff * inv) >>> 6, 0, 60)   // 有符号乘 9x13 -> 22bit
  H   = 扇区基值 + h                       // delta==0 时强制 H=0
\end{Verilog}

为什么差值必须 9 bit 有符号？因为 G-B 可以是负数（比如纯红 G=B=0 没问题，
但 (0,255,0) 的 B-R = -255）。RTL 里写成：

\begin{Verilog}
wire signed [8:0] diff_r;
always @(*) begin
    if (is_r)      diff_r = $signed({1'b0, g}) - $signed({1'b0, b}); // G-B
    else if (is_g) diff_r = $signed({1'b0, b}) - $signed({1'b0, r}); // B-R
    else           diff_r = $signed({1'b0, r}) - $signed({1'b0, g}); // R-G
end
\end{Verilog}

\begin{warnbox}{位精确的三个要点（每次改这个模块都要默念）}
\textbf{1.} Python 的 \code{>>} 对负数就是向下取整，Verilog 必须用
\code{>>>}（算术右移）配合 \code{signed} 才能一致；用 \code{>>} 或漏掉
signed 会让负半区（纯红附近）差 1 度。\\
\textbf{2.} \code{4'sd0} 这类字面量如果不显式 signed，混合运算会整体
变无符号，差值直接回绕成大正数。\\
\textbf{3.} 金标准里 \code{h\_r = clip(((g-b) * inv) >> 6, 0, 60)} 用的是
numpy 的 floor 除法；RTL 的 \code{h\_raw[6:0]} 截位前先 clamp 到 0，
两边对负数语义必须逐位对齐——用 tb\_rgb2hsv 的 10010 个向量兜底。
\end{warnbox}

\subsection{已知特性：品红为什么是 H=0}

本实现把 B 最大一律归到 240+ 扇区，且内插差值取 R-G。对标准品红
(255, 0, 255)：R-G = 255-0 = 255，inv[255]=15，h = (255x15)>>6 = 59，
H = 240+59 = 299——看起来对。但对 (128, 0, 128) 这类"暗品红"：
R-G = 128，inv[128]=30，h = (128x30)>>6 = 60，H = 300，
而标准公式给 300 附近，也在窗内。

真正会翻车的是"红和品红同时出现"的场景：品红 (255,0,255) 的 R 与 B 相等，
扇区判定 \code{is\_r = (R>=G \&\& R>=B)} 会把它判成\textbf{红色扇区}
（R 优先），差值 G-B = -255，h = (-255x15)>>6 = -60，clamp 到 0，
于是品红拿到 H=0，也就是"和纯红一模一样"。

\begin{tipbox}{这是特性不是 bug}
工程测下来：红色目标检测不受影响（H=0 本来就在红色窗口内），
测试场景里放的品红干扰物会被"红窗 + S/V 门限"部分吸收或部分误判——
方案对比文档里记录了取舍结论：对比赛演示够用，若后续需要严格区分
红与品红，需要扩到 6 扇区全内插（已有演进方案，未上马）。
详见 \code{01\_方案对比/}——那里有完整的公式推导与实测数据。
\end{tipbox}

% ============================================================
\section{模块详解：hsv\_inv\_lut（生成文件）}

这个 ROM 不允许手写：由 \code{tools/gen\_vectors.py} 生成，头注释明确
"自动生成、禁止手改"。生成逻辑就一行：

\begin{Verilog}
inv[d] = 3840 // d      // d = 1..255, 用整除（floor）
\end{Verilog}

为什么不用 \code{\$readmemh} 初始化 ROM？因为要过综合——生成器直接吐出
一个 256 项的 \code{case} 表 RTL，综合器认它是常量，会映射成 Block ROM。
Python 金标准用 \code{build\_inv\_lut()} 建同一个表，
两边同源，差异只可能来自手改——所以禁止手改。

% ============================================================
\section{模块详解：hsv\_threshold}

58 行的极简模块，但有两个设计点值得学：

\begin{Verilog}
// H 窗口环绕：h_lo > h_hi 时按"或"判断（红色 335..25）
assign in_h = (h_lo <= h_hi) ? (I_h >= h_lo) && (I_h <= h_hi)
                             : (I_h >= h_lo) || (I_h <= h_hi);
assign mask_c = in_h && (I_s >= s_lo) && (I_v >= v_lo) && (I_v <= v_hi);
\end{Verilog}

\textbf{1. 8 组预置的 case ROM}：\code{I\_preset\_idx} 由按键/拨码切换，
0 红 1 绿 2 蓝 3 橙黄 4 青 5 品红 6 白 7 黑。参数表与金标准
\code{HSV\_PRESETS} 逐项同值（评审时可以直接对表）。\\
\textbf{2. 为什么阈值不归一化}：S 用 delta（未除以 max）、V 用 max，
都是 8 bit 天然值。这样省掉一次除法，且实拍标定下来门限直观
（"红色够不够纯看 S>=80"）。桌面上评测用彩色球/色卡标定过一次参数，
写死在预置表里。

% ============================================================
\section{模块详解：bin\_smooth3x3——从 v3 到 v7 的重构}

这是 Step2 成本最高的一段，也是全工程最有价值的复盘。建议对照
\code{04\_迭代文件/迭代记录\_v1到v7.md} 一起读。

\subsection{为什么是"多数表决"而不是均值}

均值滤波会把 1 像素宽的小目标糊成灰（掩码只有 0/1）；
3x3 多数表决（9 格中 >=5 个 1 就输出 1）对孤立噪点（1 个 1 周围 8 个 0）
和孔洞（8 个 1 中 1 个 0）都能正确修复，且对目标边缘保持锐利。
测试向量里专门造了"椒盐噪声 + 小方块"的帧来验证这一点。

\subsection{v3 架构：移位延迟线（失败的开始）}

第一版按"教科书"写法：

\begin{Verilog}
// v3 思路（已废弃）
row_delay: 每行一个 1280 位移位寄存器, 输出接下一级
z1 <= in; z2 <= z1;              // 水平 2 级抽头
win = {z2, z1, in} 三列各取三行版本相加
\end{Verilog}

想得很美：数据流进移位线，右端自然"看到"过去 2 个像素、过去 2 行。
实测三连败：

\begin{warnbox}{v3 的三轮失败（迭代 5、6）}
\textbf{失败 1（首错 px=13 行 0）}：移位线读端取错——链最右端是最新数据，
代码却按最老读，窗口三列取成 \code{\{k-3, k-2, k-1\}}，整体差 2 拍。
修复：读端改 bit0，窗口改 \code{\{t1, t0, cur\}}。\\
\textbf{失败 2（行 1 起全错）}：垂直延迟线"深度 W"实际只延迟了 W-1 拍
（输入端先打一拍才进线），差 1 拍。修复：线加深到 W+1。\\
\textbf{失败 3（行 5 起每行棘轮 +1 错位）}：虚槽 + 预取架构的推入/弹出
收支不平衡——每行推入 65 个但只读出 64 个，计数器每行 +1 环绕吃数据。
\textbf{这是结构性缺陷，补丁打不下去}。
\end{warnbox}

\subsection{推倒重来：v7 索引行缓存（自对齐结构）}

复盘结论写进了迭代记录："自对齐结构优先于任何需要手工推演相位的紧凑方案。
省下的调试时间是数十倍逻辑资源的代价也值得。"

v7 的核心思想一句话：\textbf{读指针 = 写指针 = 当前列 x}。

\begin{Verilog}
(* ram_style = "distributed" *) reg line0_arr [0:W-1];   // row y-1
(* ram_style = "distributed" *) reg line1_arr [0:W-1];   // row y-2

wire d0 = line0_arr[x];          // 先读: 当前列的上两行旧值
wire d1 = line1_arr[x];
always @(posedge I_clk) begin    // 后写: 当前像素进 line0, 挤出者进 line1
    if (I_valid) begin
        line1_arr[x] <= d0;
        line0_arr[x] <= I_mask;
    end
end
\end{Verilog}

同一个 \code{x} 下，读到的 \code{d0} 就是"上一行同列"，\code{d1} 就是
"上上行同列"——垂直方向不需要任何相位推导，因为读写指针天然一致。
再配 2 拍水平 FF 抽头（\code{x-2, x-1}）+ 当拍 \code{x}，3x3 窗口齐了。
这个结构直接抄自官方例程 edge\_detect\_96（第 9 课官方代码），
"官方例程当参考"是本工程的一个方法论。

\subsection{行尾窗口：一个容易漏的输出约定}

每行最后一个像素（中心 W-1）的 3x3 窗口需要 {W-2, W-1, W-1 复制}
三列——但此时流已经走到下一行行首。v7 的处理：把整行的 9 个抽头寄存器
一直留着，在下一行 \code{x==0} 拍用"行尾窗口"公式把上一行最后一个
输出补出来：

\begin{Verilog}
wire [3:0] win_last = {3'd0, t1} + {3'd0, t0} + {3'd0, t0} +
                      {3'd0, m1} + {3'd0, m0} + {3'd0, m0} +
                      {3'd0, b1} + {3'd0, b0} + {3'd0, b0};
wire last_en = (x == 11'd0) && started;
// O_eol <= I_valid && last_en;   // 行尾窗口即每行最后一个输出
\end{Verilog}

于是输出流保证每行恰好 1280 个像素、valid 连续、eol 对齐最后一
个输出——下游 centroid\_accum 直接按流数格点即可，不用知道任何
内部相位。\textbf{这就是"输出约定比内部结构重要"的实例}。

\subsection{边界复制语义表（与金标准逐位对应）}

\begin{center}
\begin{tabular}{lll}
\toprule
\rowcolor{rowtint}\textbf{位置} & \textbf{规则} & \textbf{金标准对应} \\
\midrule
列 -1 & 复制列 0 & \code{ext[0] = row[0]} \\
列 W & 复制列 W-1 & \code{ext[w+1] = row[w-1]} \\
行 -1, -2（y=0） & 三行都取当前行 & \code{top = mid = cur} \\
行 -1, -2（y=1） & 顶中两行取 row0 & \code{top = mid = row\_mid} \\
\bottomrule
\end{tabular}
\end{center}

RTL 里对应的三行选择：\code{y==0} 时 t/m 都取当拍输入；\code{y==1} 时
t/m 都取 \code{d0}（场景 row0 被行缓存挤出后仍在 line1 位置）。

\begin{okbox}{v7 验收}
tb\_bin\_smooth：4 帧 64x48（12288 个像素）与金标准
\code{smooth3x3\_majority} 逐位一致，\textbf{12288/12288 全过}。
从此没有再出过平滑相关缺陷。
\end{okbox}

% ============================================================
\section{模块详解：centroid\_accum}

\subsection{帧累加器：三个加数器 + 一组极值}

\begin{Verilog}
wire hit = I_valid && I_mask;
// 帧首 sof 清零; 命中像素: sum_x += x; sum_y += y; cnt += 1
// bbox: seen 标志法——第一个命中像素初始化 min/max, 之后逐像素比较
\end{Verilog}

位宽账：\code{sum\_x} 最大 1280x921600 会溢出 32 位吗？不会——
sum\_x 的实际上界是 \code{count x 1279 <= 921600 x 1279 ≈ 1.18e9 < 2\^{}31}，
32 位有符号都够，取无符号更宽裕。\code{cnt} 21 位（921600 < 2\^{}20 其实
20 位够，留 1 位余量）。

\subsection{最阴的 bug：除法快照（迭代 8）}

第一版：帧末脉冲启动除法状态机，直接读 \code{sum\_x/cnt} 边算边用。
tb\_centroid 显示帧 0/1 结果错、帧 2 对。根因是\textbf{时序竞争}：

\begin{Verilog}
帧 N 的最后一个像素
  -> 除法状态机启动（需要 ~66 拍）
  -> 但下一帧的 sof 清零只要 ~9 拍就到达（流水线延迟短）
  -> 除法算到一半，sum_x/cnt 已被下一帧清零 -> 商是垃圾
\end{Verilog}

为什么帧 2 反而"对"？——因为清零点落在这帧除法算完之后的拍上，纯属巧合。

修复方案：\textbf{tick 时刻全量快照}。

\begin{Verilog}
S_IDLE: if (tick_r) begin
    snap_sumx <= sum_x; snap_sumy <= sum_y; snap_cnt <= cnt;   // 定格
    snap_minx <= bmin_x; ... snap_maxy <= bmax_y;
    if (cnt < MIN_PIXELS) st <= S_LATCH;      // 无目标, 直接锁存无效
    else begin dividend <= sum_x; ... st <= S_DIV_X; end
end
\end{Verilog}

除法与结果锁存\textbf{只读 snap\_* 寄存器}，与活动累加器完全解耦。
这个改动之后，tb\_centroid 三帧全过。教训一句话：
\textbf{跨帧/跨段使用一个累加结果前，先把它"定格"下来}。

\subsection{恢复法除法与"商末位预测"}

除法用最经典的逐位恢复法（shift-subtract），32 位商 32 拍：

\begin{Verilog}
wire [53:0] rem_sh  = {remainder, dividend[31]};   // 余数左移一位, 拉进被除数最高位
wire [21:0] rem_sub = rem_sh[21:0] - divisor;
wire        rem_ge  = (rem_sh[21:0] >= divisor);
// 够减: 余数换差, 商补 1; 不够减: 余数左移结果保留, 商补 0
quotient <= {quotient[30:0], rem_ge};
dividend <= {dividend[30:0], 1'b0};
\end{Verilog}

一个小细节：最后一个商的 bit 是在 \code{bit\_cnt==1} 那一拍才算出来的，
状态机下一拍就走 \code{S\_LATCH}，来不及用 \code{quotient}。所以
\textbf{商末位与最后一次移位同拍预测锁存}：

\begin{Verilog}
if (st == S_DIV_X && bit_cnt == 6'd1)
    q_x <= {quotient[30:0], rem_ge};      // 提前拼上最后一位
\end{Verilog}

X/Y 两轴串行复用同一个除法器（各 32 拍，共 64 拍），面积省一半。
除零保护不用专门写：\code{cnt < MIN\_PIXELS(30)} 时根本不启动除法，
直接走 S\_LATCH 锁存无效结果——\textbf{最小像素门限同时兼任了除零守卫}。

\subsection{有效性、丢帧计数与输出}

\begin{Verilog}
if (snap_cnt >= MIN_PIXELS) begin
    O_target_valid <= 1'b1; O_cx <= q_x[10:0]; O_cy <= q_y[10:0];
    O_count <= snap_cnt;  O_lost_cnt <= 16'd0;
end else begin
    O_target_valid <= 1'b0; O_cx <= 0; O_cy <= 0;
    if (O_lost_cnt != 16'hFFFF) O_lost_cnt <= O_lost_cnt + 16'd1;  // 饱和
end
\end{Verilog}

\code{lost\_cnt} 饱和计数供云台方向做"连续丢目标"策略；结果在
\code{S\_LATCH} 拍通过 \code{O\_results\_upd} 发出一个脉冲——这个脉冲
同时被用做跨域握手（Step4 的 68bit 目标包 CDC load 信号）。

% ============================================================
\section{模块详解：detect\_top（窃流旁路顶层）}

\subsection{为什么挂在 96bit 点的 AWB 之后}

官方 IPC 链路里，96bit AXI-Stream 的第一站是 ISP 流水线：曝光、白平衡、
去噪、转换。挂在 \textbf{AWB（自动白平衡）之后、data128\_96 之前}，
意味着我们拿到的是"颜色已经校正过"的像素——红色球在偏黄灯光下
也能稳定落在红窗里。如果挂在 AWB 之前，阈值就得跟着灯光漂。

\subsection{纯旁路的写法}

\begin{Verilog}
assign O_tready = 1'b1;
assign O_tdata  = I_tdata;      // 逐拍零延迟直通
assign O_tvalid = I_tvalid;
assign O_tuser  = I_tuser;
assign O_tlast  = I_tlast;
\end{Verilog}

检测链从 \code{I\_tdata} 并行分支出去，\textbf{不反压、不插拍、不修改
一比特视频}。tb\_detect\_chain 里专门有一个"直通一致性"检查：整个
1382400 拍（3 帧 x 720 x 320 beat）逐拍比对输出与输入，
任何一拍的 tdata/tvalid/tuser/tlast 不同就报 PASSTHRU\_MISM。

这个结构让"上板挂载"变成一句话的事：把视频线接进来、把线接出去，
检测结果从侧口出来——不用懂官方 ISP 内部任何细节。

\subsection{顺手带出的三路输出}

\begin{itemize}
  \item \code{O\_s\_*}：串化后的 1ppc 原始流（供 Step4 OSD 当底图）；
  \item \code{O\_fid}：8 位采集帧号，每个输入帧 sof 自增（进阶（3）延迟可视化）；
  \item \code{O\_dbg\_mask}：平滑后掩码流，仿真里可直接抓波形看"机器眼里的目标"。
\end{itemize}

% ============================================================
\section{测试平台设计}

\subsection{向量生成器：所有期望值只有一个来源}

\code{tools/gen\_vectors.py} 统一生成 5 组 TB 的激励与期望：

\begin{center}
\begin{tabular}{llrl}
\toprule
\rowcolor{rowtint}\textbf{TB} & \textbf{向量} & \textbf{规模} & \textbf{构造方法} \\
\midrule
tb\_rgb2hsv & rgb2hsv\_in/exp & 10010 & 10 个角点 + 10000 随机 RGB \\
tb\_hsv\_threshold & threshold\_in/exp & 4000 & 随机 H/S/V，8 组轮转 \\
tb\_bin\_smooth & smooth\_in/exp & 12288 & 4 帧 64x48 噪声 + 移动方块 \\
tb\_centroid & centroid\_in/exp & 27648 & 3 帧 128x72 派生几何 \\
tb\_detect\_chain & chain\_beats/flags/exp & 691200 & 3 帧全尺寸真实场景 \\
\bottomrule
\end{tabular}
\end{center}

角点必须进向量集：\code{(0,0,0)、(255,255,255)、纯三原色、纯黄青品红、
灰、暗红 (224,32,32)}——里面藏着 delta=0、差值为负、clamp 触发
三个边界，随机向量撞不上这些。

\subsection{tb\_detect\_chain：全尺寸真实场景}

三帧 1280x720 是"整链验收"：

\begin{center}
\begin{tabular}{llll}
\toprule
\rowcolor{rowtint}\textbf{帧} & \textbf{场景} & \textbf{金标准结果} & \textbf{考察点} \\
\midrule
0 & 红方块 (560,300) 90x90 + 椒盐 & count=8096, cx=604, cy=345 & 噪声不破坏质心 \\
1 & 红块 + 品红干扰物 + 噪声 & count=32592, bbox 横跨 300..939 & 干扰物部分吸收 \\
2 & 小目标 15x15 (100,101) & count=221, cx=107, cy=108 & 小目标仍过门限 \\
\bottomrule
\end{tabular}
\end{center}

注意第 0 帧：方块 90x90 = 8100 像素，平滑后 8096——边缘 4 个像素被
多数表决判成背景，这个 -4 是\textbf{金标准算出来的同一个数}，
不是近似。位精确验证就是"连误差都要对得上"。

\subsection{三个向量文件格式坑（全部来自实战）}

\begin{warnbox}{生成器与 TB 的对齐协议}
\textbf{1. hex 字段宽度 = 向量位宽 / 4，不许溢出}：
\code{chain\_flags.hex} 最初每行写 2 个十六进制字符，但 tuser/tlast 合并
只有 2 bit，\code{\$readmemh} 直接按位宽截断，两个标志全毁。改成单字符
（值 0..3）。\\
\textbf{2. 期望文件用 \$fscanf 读，字段宽度必须自洽}：
\code{chain\_exp.txt} 最初 y 坐标用 \code{\%02X} 写，y=300 写出 "12C"
三个字符，\code{\$fscanf} 按固定宽度切字段全部错位。规范：坐标类字段
一律 \code{\%03X}。\\
\textbf{3. 生成器坐标不许越界}：质心向量生成器最初把方块写到画幅外，
numpy 静默裁剪后金标准少算了一块像素，TB"错得有理说不清"。
生成器末尾必须断言所有几何在画幅内。
\end{warnbox}

% ============================================================
\section{验证结果}

\subsection{回归矩阵}

\begin{center}
\begin{tabular}{llll}
\toprule
\rowcolor{rowtint}\textbf{测试} & \textbf{比对粒度} & \textbf{向量数} & \textbf{结果} \\
\midrule
tb\_rgb2hsv & H/S/V 三字段逐位 & 10010 & PASS \\
tb\_hsv\_threshold & mask 逐位 & 4000 & PASS（首过） \\
tb\_bin\_smooth & mask 逐位 + 行对齐 & 12288 & PASS（v7 重构后） \\
tb\_centroid & 8 字段逐字段 & 3 帧 & PASS（快照修复后） \\
tb\_detect\_chain & 帧结果 + 直通逐拍 & 3 帧全尺寸 & PASS \\
\bottomrule
\end{tabular}
\end{center}

\code{rgb2hsv} 与 \code{hsv\_threshold} 是\textbf{一次通过}（首过）的
模块——因为它们的算术在 Python 里已经调通，RTL 是抄写而不是设计。
这正是"金标准先行"的价值：Step2 全部真缺陷都出在\textbf{结构}
（平滑相位、除法竞争）而不是算术上。

\subsection{Step2 阶段抓到的真缺陷}

\begin{center}
\begin{tabular}{llp{5.2cm}l}
\toprule
\rowcolor{rowtint}\textbf{编号} & \textbf{现象} & \textbf{根因} & \textbf{类型} \\
\midrule
S1 & 平滑行 0 首错 & 移位线读端相位取反 & RTL 结构 \\
S2 & 平滑行 1 起错 & 延迟线实际深度 W-1 & RTL 结构 \\
S3 & 每行棘轮 +1 & 虚槽架构收支不平衡 & RTL 架构 \\
S4 & 质心帧 0/1 错 & 除法与下帧清零竞争 & RTL 时序 \\
S5 & 整链编译错 & sim\_video\_src 缺 timescale & TB \\
S6 & 标志全毁 & flags hex 超位宽截断 & TB 脚本 \\
S7 & 字段错位 & 期望文件字段宽度溢出 & TB 脚本 \\
\bottomrule
\end{tabular}
\end{center}

S3 是最贵的一课（v3 到 v6 三个补丁都没救活），最后靠 v7 自对齐结构
一次性解决；S4 教会"结果先定格"。两个教训都写进了迭代记录供
后面 Step3/Step4 复用。

\begin{okbox}{Step2 验收口径}
1）5 个 TB 全绿；2）整链 921600 像素直通逐拍一致；3）三帧真实场景
结果与 Python 金标准逐字段一致（含 count 的每个整数）。达到这三条，
Step2 视作"无硬件条件下的位精确完成"。
\end{okbox}

% ============================================================
\section{上板挂载说明（给上板组）}

\begin{center}
\begin{tabular}{ll}
\toprule
\rowcolor{rowtint}\textbf{项目} & \textbf{内容} \\
\midrule
挂载点 & 官方 isp\_top 输出 96bit AXI-Stream（AWB 后、data128\_96 前） \\
本例顶层 & \code{detect\_top}（纯旁路，O\_tdata 等直接回接官方链路） \\
时钟 & 与视频像素时钟同域（74.25 MHz），无需 CDC \\
复位 & 与官方链路同复位（低有效） \\
按键 & \code{I\_preset\_idx[2:0]} 接 KEY 或拨码（换目标颜色） \\
输出 & 68bit 目标包 + results\_upd 脉冲（送 Step3 云台 / Step4 OSD） \\
资源估计 & 行缓存 2x1280 bit（dis-RAM）+ 256x12 ROM + 1 个 9x13 乘法 + \\
         & 32bit 除法器——在 PH1P35 上占比很小（TD 综合报告见 Step5） \\
\bottomrule
\end{tabular}
\end{center}

上板三步：
\begin{enumerate}
  \item 在 TD 工程里把 6 个 RTL 文件加入源列表（注意 \code{hsv\_inv\_lut.v}
        是生成文件，别手改）；
  \item 实例化 \code{detect\_top}，把官方 96bit 流接 I\_*、输出接回；
  \item 用预定按键做 \code{I\_preset\_idx}，输出挂 ILA/串口观测（选做）。
\end{enumerate}

\begin{tipbox}{和 Step1 的分工}
Step1（视频链路打通）已经把串化器、CDC、测试图案源验证过了；
Step2 的 detect\_top 里直接实例化 Step1 的 ser96to24。两个阶段在
仓库里是独立目录、独立 TB，但上板时合成一个顶层——
接口文档（\code{06\_interfaces}）是它们之间的合同。
\end{tipbox}

% ============================================================
\section{小结：Step2 的可复用产出}

\begin{itemize}
  \item \textbf{6 个可综合 RTL 模块}（916 行），全部通过位精确验证；
  \item \textbf{一份金标准} \code{golden\_model.py}（被 Step3/4/5 继续复用）；
  \item \textbf{5 组向量 + 5 个 TB}，可以随时一键回归；
  \item \textbf{一条方法论}：金标准先行、自对齐结构、结果先定格、
    标志随行——后面的 Step3/4/5 全部照此执行；
  \item \textbf{一份踩坑档案}（\code{04\_迭代文件}），新队员接手时
    先读它，能省下至少一周的重复调试。
\end{itemize}

\begin{okbox}{下一步（Step3 预告）}
Step2 的输出是"目标坐标"；Step3 要把它变成云台的\textbf{运动}：
误差计算、Q8.8 定点比例控制、舵机 PWM——以及一次真实的振荡事故
（Kp=90 回路增益 1.3 导致跑飞 109 度）的完整复盘。
\end{okbox}


% ============================================================
\clearpage
\part*{结语}
\section*{总结与下一步}
\addcontentsline{toc}{section}{总结与下一步}

\subsection*{Step2 学到了什么}
\addcontentsline{toc}{subsection}{Step2 学到了什么}
\begin{itemize}
  \item 颜色空间变换：RGB 判颜色为什么不稳，HSV 怎么解决
  \item 定点算术：查表代替除法、移位代替乘除、floor 与 signed 的坑
  \item 阈值分割：H 窗口环绕（红色跨 0 度）与 S/V 门限的配合
  \item 3x3 多数表决去噪：行缓存 + 窗口 + 边界复制
  \item 质心/包围框：帧累加、快照除法、最小像素门限
  \item 位精确验证方法学：金标准先行、向量生成、逐位比对
\end{itemize}

\subsection*{下一步学什么}
\addcontentsline{toc}{subsection}{下一步学什么}
\begin{itemize}
  \item \textbf{Step3 云台闭环}：把坐标变成运动
  \begin{itemize}
    \item 误差计算与死区（目标在画面中心小范围不动）
    \item Q8.8 定点比例控制（KP 增益整定，回路增益 < 1 才不振荡）
    \item 舵机 PWM 生成与单帧限速
    \item 振荡事故复盘：Kp=90 为什么跑飞 109 度
  \end{itemize}
  \item \textbf{Step4 OSD 与整链}：把结果画回画面（文字/框/延迟可视化）
  \item 对应文档：\code{03\_vision/step3\_closed\_loop/}、\code{step4\_osd\_integration/}
  \item 对应金标准：\code{tools/golden\_model.py::gimbal\_ctrl\_step}
\end{itemize}

\end{document}
